\section{CPU shape construction: \code{Octree::apply}}
\label{sec:cpushape}

\subsection{Entry point and root expansion}

An edit is expressed as
\begin{lstlisting}[caption={The public shape entry point (simplified).},label={lst:apply}]
void Octree::apply(operation, function, model, painter, coordPainter,
                   minSize, simplifier, updateHandler, deleteHandler) {
    unique_lock treeLock(treeMutex);
    ShapeArgs args = ...;       // operation, function, model, painter, simplifier
    expand(args);               // grow the root if the new shape lies outside
    configureShapeCache(args);  // lattice origin/spacing for this apply
    OctreeNodeFrame frame(root, nullptr, *this, ...);
    NodeOperationResult r;
    shape(r, frame, args, &localContext, updateHandler, deleteHandler);
}
\end{lstlisting}

\code{expand} grows the root by powers of two until the new primitive is
fully contained.  It computes the child index of the primitive's centre in
the current root cube and places the old root as that child of a new,
doubled root.  This is the only operation that changes the world cube; the
existing subtree is preserved bit-for-bit.  The lattice origin and leaf
length used by the SDF cache are computed \emph{after} expansion.

\subsection{The recursive \code{shape} call}

For a cube of length $L$ and a target leaf size $m$, the recursion has
three regimes:
\begin{enumerate}
\item \textbf{Leaf:} $L\le m$ and the node is a leaf.  Sample the eight
  corners of the shape field, combine them with the existing field, and
  classify the result.
\item \textbf{Interior with early out:} the node is not a leaf, but the
  shape field's centre value and bounding-volume test prove that the
  primitive cannot affect this cube.  Copy the existing field and stop.
\item \textbf{Interior descend:} recurse into the eight children, join the
  results, classify the parent from the children, and optionally simplify.
\end{enumerate}

\begin{figure}[htbp]
\centering
\begin{tikzpicture}[node distance=6mm and 8mm,font=\small]
  \node[box] (enter) {enter \code{shape}(cube, frame)};
  \node[box,below=of enter] (leaf) {$L\le m$ and node leaf?};
  \node[box,right=of leaf] (sample) {sample 8 shape corners\\combine with existing field\\classify result};
  \node[box,below=of leaf] (outside) {shape centre outside\\expanded cube?};
  \node[box,right=of outside] (copy) {copy existing field\\no child work};
  \node[box,below=of outside] (children) {recurse 8 children\\(parallel fork band)};
  \node[box,below=of children] (parent) {parent SDF from children\\classify result};
  \node[box,below=of parent] (simp) {simplify?\\set \code{selectedLod}};
  \node[box,below=of simp] (paint) {paint leaf / dominant child\\write node fields};
  \node[box,below=of paint] (publish) {if \code{chunkLod>0} and Surface:\\bump version, call update handler};
  \draw[arrow] (enter)--(leaf);
  \draw[arrow] (leaf)--node[right,note]{yes}(sample);
  \draw[arrow] (leaf)--node[left,note]{no}(outside);
  \draw[arrow] (outside)--node[right,note]{yes}(copy);
  \draw[arrow] (outside)--node[left,note]{no}(children);
  \draw[arrow] (children)--(parent);
  \draw[arrow] (parent)--(simp);
  \draw[arrow] (simp)--(paint);
  \draw[arrow] (paint)--(publish);
  \draw[arrow,dashed] (sample) |- (paint);
  \draw[arrow,dashed] (copy) |- (publish);
\end{tikzpicture}
\caption{The shape recursion.  A leaf samples the new field; an unaffected
interior copies the old field; an affected interior recurses, classifies
the parent from its children, simplifies, paints, and (for surface chunk
roots) publishes a new version.}
\label{fig:shapeflow}
\end{figure}

\subsection{Sampling and caching}

The shape field is sampled at cube corners with \code{evaluateSDF}.  A
naive implementation calls the primitive's distance function once per
query.  The engine uses a per-apply cache keyed on the apply's leaf lattice:
\begin{equation}
  \mathbf{x} = \mathbf{o} + \mathbf{n}\,h,
\end{equation}
where $\mathbf{o}$ is the lattice origin and $h$ the half leaf length.
An integer lattice key is accepted only when the reconstructed float is
bit-identical to the query.  Anything off-lattice (a stored tree from an
earlier apply with a smaller leaf size) falls back to a float-keyed map.
The cache is a pure optimisation: the returned value is the value the
primitive would produce.  The cache is per-thread, generation-tagged, and
reset when the generation changes.  This avoids both a global lock and
stale values across applies.

\subsection{The three early-out rules}

The recursion has three important early outs.  They are necessary for
performance and for correctness at the field's infinity regions.

\paragraph{All-infinity node.}
If no corner of the existing field has data, the node is untouched.  The
operation may propagate from infinity (a union with a finite shape) or not
(a subtraction whose negative operand is undefined).  If the shape's
centre is more than the cube's half-diagonal away from the solid, the
whole cube is classified Solid or Empty without descending.

\paragraph{Solid-preserving operation.}
If the existing cube is Solid, all its corners finite, and the operation
preserves solidity, the result is still Solid.  The operation only needs
to combine the existing and shape corner values; no child descent is
needed.  This is the common case for a union of overlapping solids.

\paragraph{Solid cube and centre test.}
If the existing cube is Solid and the combined centre is more than the
cube's half-diagonal inside the surface, the whole cube is Solid.  This
avoids descending into regions the new shape cannot reach.  The test uses
the half diagonal $L\sqrt{3}/2$, not the half edge: a sphere can enter
through a corner.

If none of the early outs applies, the function's \code{check(cube)} test
is consulted.  A \code{Disjoint} result skips the children; a
\code{Contains} result can be handled by the same solid/empty rules.
Finally, if the node is a new surface node with no existing children, the
recursion is forced to descend even if the centre is outside, because the
surface may still cross a corner.

\subsection{Result composition}

For a non-leaf node, each child returns eight result SDF values.  The
parent's result array is the concatenation of the children's corner arrays
in child order.  The parent's type is:
\begin{equation}
  \operatorname{parentType} =
  \begin{cases}
    \text{Empty}, & \text{all children Empty},\\
    \text{Solid}, & \text{all children Solid},\\
    \text{Surface}, & \text{otherwise}.
  \end{cases}
\end{equation}
This is exact because the child corner arrays share the parent's corners
on their common faces.  The parent does not need to re-evaluate the field.

\subsection{Painting}

Painting happens at the shape leaf, where the node's material is decided.
Brushes can be constant (a material index), procedural (terrain
classification from height and slope), or position-dependent (the
\code{LandBrush} classifier).  The painter writes:
\begin{itemize}
\item \code{brushIndex} (material),
\item \code{hsv} (tint),
\item \code{textureType} (mapping technique),
\item \code{data} (UV scale or authored UV).
\end{itemize}
Internal nodes inherit the dominant painted child's state.  The exact
dominance rule is intentionally simple: the last surface child with a
valid brush that wins the per-child tally sets the parent.  The tally is
kept per child (not across all eight) so the historical result is
preserved; changing it would change the mesh's material assignment and UV
scale.

\subsection{Node publishing and versions}

A node that is both a surface and a chunk root (\code{chunkLod>0}) is a
published chunk.  At the end of its recursion, the engine:
\begin{enumerate}
\item sets its type, SDF, chunk flag, levels, brush, mapping state;
\item increments \code{version};
\item calls the update handler with \code{(level, node, cube)}.
\end{enumerate}
The handler queues a \code{PendingMeshData}.  The renderer later flattens
the neighbourhood, dispatches the GPU walk, and publishes the span.  A
version is the only identity the renderer trusts: the same node address may
be reused by the allocator, so pointer identity alone is not a revision.

\subsection{Threading}

The recursion forks when the child cube length lies in
\begin{equation}
  \left(\frac{m\cdot 16}{4},\;m\cdot 16\right],
\end{equation}
where $m$ is the target leaf size.  A forked node enqueues eight tasks into
the octree's thread pool; each task owns a copy of its frame and the shape
arguments and writes into a shared result block.  The parent joins all
eight futures before reading the block.  The fork band is two levels deep:
a $16m$ node forks eight $8m$ children and each of those forks eight $4m$
grandchildren; the $4m$ nodes do not fork again.  This bounds the number of
tasks per fork to 64 and avoids a thundering herd at the leaf level.

The join loop is careful about pool starvation.  Instead of blocking on a
future while the pool is idle, it calls \code{runOneTask()} to execute
queued work on the waiting thread; only when the queue is empty does it
block.  This is what keeps a nested fork from deadlocking on a pool with
fewer threads than the fan-out.

\subsection{Why the CPU never solves a vertex}

The current shape path deliberately avoids \code{SDF::getPosition} and
\code{getNormalFromPosition}.  The only host callers are the offline
Blender exporter.  This is not just a performance choice: it is a
correctness boundary.  If the runtime computed a vertex on the host and the
GPU computed another, the two would differ in the last bits and the
renderer's mesh would not match the exporter's.  By making the GPU the sole
Cartesian producer, the runtime has one answer and the exporter has a
reference implementation that can be compared offline.
